%!TEX root = ../thesis.tex
%*******************************************************************************
%*********************************** First Chapter *****************************
%*******************************************************************************

\chapter{Voter Candidate Range Profile Detection}  %Title of the First Chapter

\ifpdf
    \graphicspath{{Chapter4/Figs/Raster/}{Chapter4/Figs/PDF/}{Chapter4/Figs/}}
\else
    \graphicspath{{Chapter4/Figs/Vector/}{Chapter4/Figs/}}
\fi

% ========================================================================
% VCR Profile Detection Algorithm
% ========================================================================

This chapter describes in detail and compares created algorithms for the VCR property detection.

\section{Overview}
Detection algorithms are used to check whether election profile satisfies desirable structural properties.
For example to check CR (VR) properties it's enough to check for the \textit{Consecutive Ones Property} in the preferences matrix.

My focus lies on the one dimmensional VCR property. 
Having preferences matrix as input we have to provide mapping of each agent (voter/candidate) to the point on real line
and nonnegative number representing radius, so that it satisfies the VCR property.  

\section{Integer Linear Programming}
ILP is a common approach to solve NP-hard problems, 
however it still can be used effictively for other problems.
First we have to define the problem as an integer linear program.
(WIP: Czy w ogole musze tlumaczyc co to ILP?) 
That is a system of linear inequalities and constraints and 
optionally linear function to optimize, alterinatively we can just check feasibility of the constraints.
To obtain a solution we have to use a optimization solver.
Popular choice are Gurobi or CPLEX which are state of the art general purpose solvers.

Solution based on the ILP is excellent choice for the first algorithm to develop.
It provides efficient and relatively simple detection algorithm,
Additionally it can be used as a benchmark for checking future solutions.

\subsubsection{Transform VCR definition to linear constraints}
Before forumulating ILP we have to redefine the VCR definition beacasue
it is based on the euclidean distance hence contains quadratic function.
Equation \ref{vcr-to-linear-eq} shows transforamtion of VCR equation into 
linear inequalities.

\begin{equation}\label{vcr-to-linear-eq}
    \begin{split}
    VCR & \iff \sqrt{(x_c - x_v)^2} \le r_c + r_v \\
     & \iff \|x_c - x_v\| \le r_c + r_v \\
     & \iff -(r_c + r_v) \le x_c - x_v \le r_c + r_v \\
     & \iff x_c - x_v \le r_c + r_v\ \ \lor \ \ -(x_c - x_v) \le r_c + r_v
    \end{split}
\end{equation}
The constraints obtained from equation \ref{vcr-to-linear-eq} 
still are not suitable for ILP beacasue it contains $\lor$ logical condition
which cannot be directly represented in the ILP.
However this issue is quite common and most solvers provide additional mechanisms to solve it.

In Gurobi solver it's called \textit{INDICATOR} constraint
which allows to create versatile constraints like:
$z = f \rightarrow a^Tx \leq b$ which states that if the binary indicator variable $z$
is equal to $f$, where $f \in \{0,1\}$, then the linear constraint 
$a^Tx \leq b$ should hold. On the other hand, if $z = 1-f$, 
the linear constraint may be violated.
This approach is losely based (WIP Fact check) on the so called \textit{BigM} method which is low level solution, and requires parameter tunning depending on the problem.

\subsubsection{Detection algorithm}

Introduce new binary variable for each pair of candidate and voter:\\
$z_{cv}$


\begin{alignat*}{2}
    & \text{minimize: } & & \sum_{j=1}^{m} w_{j}x_{j} \\
     & \text{subject to: }& \quad & \sum_{\mathclap{{j:e_{i} \in S_{j}}}}\begin{aligned}[t]
                      x_{j} & \geq 1,& i & =1, \dots, n\\[3ex]
                    x_{j} & \in \{0,1\}, & \quad j &=1 ,\dots, m
                  \end{aligned}
\end{alignat*}

\begin{align*}
\textit{Minimize 0} \\ 
\textit{foreach (c,v) in } C \times V : \\\textit{if } c \in ApprovalSet(v) \textit{ add constraints} \\ 
x_c - x_v \le r_c + r_v \\
-(x_c - x_v) \le r_c + r_v \\
\textit{if } c \not\in ApprovalSet(v) \textit{ add constraints} \\
WIP Replace with IndicatorConstraint \\
x_c - x_v  + M(1 - z1_{cv}) \ge r_c + r_v + 1 \\
-(x_c - x_v) + M(1 - z2_{cv}) \ge r_c + r_v + 1 \\
z1_{cv} + z2_{cv} >= 1
\end{align*}

\section{(WIP) My VCR detection algorithm}

\section{Performance Comparison}
How fast Gurobi Solver can be?